\usepackage[utf8]{inputenc}
\usepackage[T1,T2A]{fontenc}
\usepackage[russian]{babel}
\usepackage{color}
\usepackage{calc}
\usepackage{graphicx}
\usepackage{epstopdf}
\usepackage{hyperref}
\hypersetup{unicode,colorlinks}
\usepackage{csquotes}
\usepackage{upquote}
\usepackage{cprotect}
\usetheme[progressbar=head,numbering=fraction,block=fill]{metropolis}
\usepackage{dejavu}
\usepackage{etoolbox}
\usepackage{bm}
\usepackage[ND]{prftree}
\usepackage[tableaux]{prooftrees}
\usepackage{mathtools} % for bigtimes
\usepackage{diagbox}
\usepackage{verbatimbox}
\usepackage{fitch}
% fitch по умолчанию печатает эти правила как $\Rightarrow$I/E; в курсе импликация всюду $\to$.
\gappto\ndrules{\def\ii{\by{$\to$I}}\def\ie{\by{$\to$E}}}
%\usepackage[defaultsans]{droidsans}
%\usepackage[defaultmono]{droidmono}
%\makeatletter
%\input droid/t1fdm.fd
%\makeatother

\makeatletter
\DeclareRobustCommand*{\bfseries}{%
    \not@math@alphabet\bfseries\mathbf
    \fontseries\bfdefault\selectfont
    \boldmath
}
\makeatother

\makeatletter
\chardef\straightquote@code=\catcode`'
\chardef\backquote@code=\catcode``
\catcode`'=\active \catcode``=\active
\patchcmd{\@noligs}
{\textasciigrave}
{\fixedtextasciigrave}
{}{}
\newcommand{\fixedtextasciigrave}{%
    \makebox[.5em]{\fontencoding{TS1}\fontfamily{fvs}\selectfont\textasciigrave}% Vera Sans
}
\catcode\lq\'=\straightquote@code
\catcode\lq\`=\backquote@code
\makeatother

\setbeamercovered{again covered={\opaqueness<1->{25}}}

\forestset{close with=\ensuremath{\bigtimes}}

\newcommand{\customframefont}[1]{
    \setbeamertemplate{itemize/enumerate body begin}{#1}
    \setbeamertemplate{itemize/enumerate subbody begin}{#1}
}

\NewEnviron{framefont}[1]{
    \customframefont{#1} % for itemize/enumerate
    {#1 % For the text outside itemize/enumerate
        \BODY
    }
    \customframefont{\normalsize}
}


\makeatletter
\DeclareRobustCommand{\rvdots}{%
    \vphantom{\int\limits^x}\smash[t]{\vdots}
}
\makeatother

\newcommand{\N}{\mathbb{N}}
\newcommand{\Z}{\mathbb{Z}}
\newcommand{\Q}{\mathbb{Q}}
\newcommand{\R}{\mathbb{R}}
\newcommand{\pow}[1]{\mathcal{P}({#1})}
\newcommand{\continuum}{\mathfrak{c}}

% \reveal<спец>{...}: показать содержимое с указанного слайда, до этого оставив пустое место той же ширины.
% В отличие от \pause и \uncover не оставляет невидимых PGF-специалов; внутри выключной формулы \pause
% оставляет их в пустом абзаце после формулы, и TeX набирает лишнюю пустую строку.
% Следующий \item после такой формулы получает <.(1)-> (тот же слайд, что и \reveal<+(1)->)
% или <+(1)-> (следующий слайд); в beamer «.» "--- это значение счётчика минус один.
\newcommand<>{\reveal}[1]{\alt#2{#1}{\phantom{#1}}}

% Многобуквенные имена в формулах: курсив с нормальным кернингом (а не произведение букв).
\newcommand{\ident}[1]{\mathit{#1}}
\newcommand{\Var}{\ident{Var}}
\newcommand{\Prop}{\ident{Prop}}
\newcommand{\Const}{\ident{Const}}
\newcommand{\Fun}{\ident{Fun}}
\newcommand{\Pred}{\ident{Pred}}
\newcommand{\Term}{\ident{Term}}
\newcommand{\Frm}{\ident{Form}} % \Form занят hyperref
\newcommand{\arity}{\ident{arity}}
\newcommand{\Th}{\ident{Th}}
\newcommand{\Dom}{\ident{Dom}}
\newcommand{\Ran}{\ident{Ran}}
\newcommand{\id}{\ident{id}}
\newcommand{\Ord}{\ident{Ord}}
\newcommand{\RAA}{\ident{RAA}}
% Названия теорий и аксиом: прямой шрифт.
\newcommand{\ZF}{\mathrm{ZF}}
\newcommand{\ZFC}{\mathrm{ZFC}}
\newcommand{\AC}{\mathrm{AC}}
\newcommand{\DC}{\mathrm{DC}}
\newcommand{\CH}{\mathrm{CH}}
\newcommand{\GCH}{\mathrm{GCH}}

\author{Алексей Романов}
\subtitle{Математическая логика и теория алгоритмов}
%\logo{}
\institute{МИЭТ}
\subject{Математическая логика и теория алгоритмов}
%\setbeamercovered{transparent}
%\setbeamertemplate{navigation symbols}{}
